% figures/w06-rotation.svg — 회전행렬 R(psi) 가 무엇인가
%
%   bash _tools/tikz2svg.sh figures/src/w06-rotation.tex figures/w06-rotation.svg
%
% 사용자 지시 (2026-10-26): "회전행렬 설명도 최대한 쉽게" — 캡스톤디자인의
% assets/src/w06-rotation.tex 를 영어로 옮긴 것이다. 구조(패널 셋·식 띠)는 그대로
% 두고, 색만 이 볼트의 팔레트로 바꾸었다. kotex 은 쓰지 않는다.
%
% 이 그림이 말하는 것 넷
%   1) 선체 축은 NED 축을 psi 만큼 돌린 것이다. 그것이 전부다
%   2) 벡터는 하나인데 어느 축에서 읽느냐에 따라 숫자 두 쌍이 나온다
%   3) R 의 **열**은 선체 축의 단위벡터를 NED 에서 읽은 것이다 — 외울 것이 없다
%   4) 거꾸로 읽는 것은 전치다. R^-1 = R^T (직교행렬)
%
% 출처 — 제어공학특론 "1. 동역학 모델" (Fossen TTK4190 / QUT Robot Academy).
%        그 자료의 rot2 예제 수치(30도)를 그대로 쓴다
%
% 색 대응 (캡스톤 -> 이 볼트)
%   guid  #1D4ED8 -> bodyc  (국소 정의. w05-guidance-loop 의 knownc 와 같은 파랑)
%   nav   #7C3AED -> accent (같은 색이다)
%   ctrl  #EA580C -> errc   #C2410C
%   thr   #D97706 -> alng   #B45309
%   plant #16A34A -> plantc (국소 정의. 볼트 팔레트에 초록이 없다)
%
% 회전 방향 확인 — psiarc 는 arc(90:60) 이고 끝 각이 작으므로 **시계방향**이다.
% 화면은 x_n(북) 위 · y_n(동) 오른쪽이므로 psi > 0 은 선수를 우현으로, 곧 화면에서
% 시계방향으로 돌린다. 맞다 (CLAUDE.md §4-6-2)
%
% 좌표 메모 (mm)
%   패널 원점  (1) 0   (2) 76   (3) 152
%   제목 y=60 · 설명 y=55 · 축 끝 y=36 · 원점 y=0 · 변수표 y=-20 · 가로줄 y=-36

\documentclass[border=6pt]{standalone}
% gnc-style.tex — 이 볼트의 TikZ 그림이 공유하는 스타일.
%
% 저널 그림 규약을 스타일 하나로 굳혀 둔다. 그림마다 선 굵기나 화살촉을
% 다시 정하지 않는다 — 그것이 그림이 제각각으로 보이는 원인이다.
%
% 쓰는 법:  % gnc-style.tex — 이 볼트의 TikZ 그림이 공유하는 스타일.
%
% 저널 그림 규약을 스타일 하나로 굳혀 둔다. 그림마다 선 굵기나 화살촉을
% 다시 정하지 않는다 — 그것이 그림이 제각각으로 보이는 원인이다.
%
% 쓰는 법:  % gnc-style.tex — 이 볼트의 TikZ 그림이 공유하는 스타일.
%
% 저널 그림 규약을 스타일 하나로 굳혀 둔다. 그림마다 선 굵기나 화살촉을
% 다시 정하지 않는다 — 그것이 그림이 제각각으로 보이는 원인이다.
%
% 쓰는 법:  \input{gnc-style}  (figures/src/ 안에서)

\usepackage{tikz}
\usepackage{amsmath}
\usepackage{xcolor}
\usetikzlibrary{arrows.meta, positioning, calc, fit, backgrounds, decorations.pathmorphing}

% ---- 색 --------------------------------------------------------------------
% 색은 강조 하나에만 쓴다. 나머지는 검정.
\definecolor{ink}{HTML}{1A1A1A}
\definecolor{accent}{HTML}{7C3AED}
\definecolor{warn}{HTML}{B91C1C}
\definecolor{soft}{HTML}{666666}

% 기하 그림(좌표계·유도)에서만 쓰는 넷. 블록선도는 위의 흑백 + accent 로 충분하다.
\definecolor{pathc}{HTML}{0D9488}   % 계획 경로 · 항적
\definecolor{errc}{HTML}{C2410C}    % 오차 (cross-track 등)
\definecolor{alng}{HTML}{B45309}    % along-track
\definecolor{flow}{HTML}{0369A1}    % 조류 · 유속

% ---- 선과 화살촉 -----------------------------------------------------------
\tikzset{
  % 신호선. 화살촉은 작고 뾰족한 것 하나뿐이다.
  sig/.style   = {ink, line width=0.5pt, -{Stealth[length=4pt, width=3pt]}},
  wire/.style  = {ink, line width=0.5pt},
  % 강조 경로 — 파선으로 구분한다. 흑백 인쇄에서도 살아야 하기 때문이다.
  asig/.style  = {accent, line width=0.5pt, densely dashed,
                  -{Stealth[length=4pt, width=3pt]}},
  awire/.style = {accent, line width=0.5pt, densely dashed},
  % 블록. 흰 바탕, 직각, 검은 테두리.
  blk/.style   = {draw=ink, line width=0.55pt, fill=white,
                  minimum width=17mm, minimum height=8mm, inner sep=2pt, font=\small},
  % 합산점
  sum/.style   = {draw=ink, line width=0.55pt, fill=white, circle,
                  inner sep=0pt, minimum size=4.6mm, font=\scriptsize},
  asum/.style  = {draw=accent, line width=0.55pt, fill=white, circle,
                  inner sep=0pt, minimum size=4.6mm, font=\scriptsize},
  % 분기점
  dot/.style   = {ink, fill=ink, circle, inner sep=0pt, minimum size=1.5mm},
  % 신호 이름 — 선 위에 작게
  sname/.style = {font=\small, inner sep=1.5pt},
  % 그림 안의 짧은 설명. 그림 안의 글은 최소로 둔다
  note/.style  = {font=\scriptsize, soft, inner sep=1.5pt},
  anote/.style = {font=\scriptsize, accent, inner sep=1.5pt},
  % 축
  axis/.style  = {ink, line width=0.5pt},
  tick/.style  = {font=\scriptsize, soft},
}

% 캡션 한 줄 — 그림 아래 가는 선과 함께
\newcommand{\gnccaption}[2]{%
  \node[anchor=north west, text width=#1, align=left, font=\scriptsize, ink]
        at (current bounding box.south west) {\vspace{2pt}#2};}
  (figures/src/ 안에서)

\usepackage{tikz}
\usepackage{amsmath}
\usepackage{xcolor}
\usetikzlibrary{arrows.meta, positioning, calc, fit, backgrounds, decorations.pathmorphing}

% ---- 색 --------------------------------------------------------------------
% 색은 강조 하나에만 쓴다. 나머지는 검정.
\definecolor{ink}{HTML}{1A1A1A}
\definecolor{accent}{HTML}{7C3AED}
\definecolor{warn}{HTML}{B91C1C}
\definecolor{soft}{HTML}{666666}

% 기하 그림(좌표계·유도)에서만 쓰는 넷. 블록선도는 위의 흑백 + accent 로 충분하다.
\definecolor{pathc}{HTML}{0D9488}   % 계획 경로 · 항적
\definecolor{errc}{HTML}{C2410C}    % 오차 (cross-track 등)
\definecolor{alng}{HTML}{B45309}    % along-track
\definecolor{flow}{HTML}{0369A1}    % 조류 · 유속

% ---- 선과 화살촉 -----------------------------------------------------------
\tikzset{
  % 신호선. 화살촉은 작고 뾰족한 것 하나뿐이다.
  sig/.style   = {ink, line width=0.5pt, -{Stealth[length=4pt, width=3pt]}},
  wire/.style  = {ink, line width=0.5pt},
  % 강조 경로 — 파선으로 구분한다. 흑백 인쇄에서도 살아야 하기 때문이다.
  asig/.style  = {accent, line width=0.5pt, densely dashed,
                  -{Stealth[length=4pt, width=3pt]}},
  awire/.style = {accent, line width=0.5pt, densely dashed},
  % 블록. 흰 바탕, 직각, 검은 테두리.
  blk/.style   = {draw=ink, line width=0.55pt, fill=white,
                  minimum width=17mm, minimum height=8mm, inner sep=2pt, font=\small},
  % 합산점
  sum/.style   = {draw=ink, line width=0.55pt, fill=white, circle,
                  inner sep=0pt, minimum size=4.6mm, font=\scriptsize},
  asum/.style  = {draw=accent, line width=0.55pt, fill=white, circle,
                  inner sep=0pt, minimum size=4.6mm, font=\scriptsize},
  % 분기점
  dot/.style   = {ink, fill=ink, circle, inner sep=0pt, minimum size=1.5mm},
  % 신호 이름 — 선 위에 작게
  sname/.style = {font=\small, inner sep=1.5pt},
  % 그림 안의 짧은 설명. 그림 안의 글은 최소로 둔다
  note/.style  = {font=\scriptsize, soft, inner sep=1.5pt},
  anote/.style = {font=\scriptsize, accent, inner sep=1.5pt},
  % 축
  axis/.style  = {ink, line width=0.5pt},
  tick/.style  = {font=\scriptsize, soft},
}

% 캡션 한 줄 — 그림 아래 가는 선과 함께
\newcommand{\gnccaption}[2]{%
  \node[anchor=north west, text width=#1, align=left, font=\scriptsize, ink]
        at (current bounding box.south west) {\vspace{2pt}#2};}
  (figures/src/ 안에서)

\usepackage{tikz}
\usepackage{amsmath}
\usepackage{xcolor}
\usetikzlibrary{arrows.meta, positioning, calc, fit, backgrounds, decorations.pathmorphing}

% ---- 색 --------------------------------------------------------------------
% 색은 강조 하나에만 쓴다. 나머지는 검정.
\definecolor{ink}{HTML}{1A1A1A}
\definecolor{accent}{HTML}{7C3AED}
\definecolor{warn}{HTML}{B91C1C}
\definecolor{soft}{HTML}{666666}

% 기하 그림(좌표계·유도)에서만 쓰는 넷. 블록선도는 위의 흑백 + accent 로 충분하다.
\definecolor{pathc}{HTML}{0D9488}   % 계획 경로 · 항적
\definecolor{errc}{HTML}{C2410C}    % 오차 (cross-track 등)
\definecolor{alng}{HTML}{B45309}    % along-track
\definecolor{flow}{HTML}{0369A1}    % 조류 · 유속

% ---- 선과 화살촉 -----------------------------------------------------------
\tikzset{
  % 신호선. 화살촉은 작고 뾰족한 것 하나뿐이다.
  sig/.style   = {ink, line width=0.5pt, -{Stealth[length=4pt, width=3pt]}},
  wire/.style  = {ink, line width=0.5pt},
  % 강조 경로 — 파선으로 구분한다. 흑백 인쇄에서도 살아야 하기 때문이다.
  asig/.style  = {accent, line width=0.5pt, densely dashed,
                  -{Stealth[length=4pt, width=3pt]}},
  awire/.style = {accent, line width=0.5pt, densely dashed},
  % 블록. 흰 바탕, 직각, 검은 테두리.
  blk/.style   = {draw=ink, line width=0.55pt, fill=white,
                  minimum width=17mm, minimum height=8mm, inner sep=2pt, font=\small},
  % 합산점
  sum/.style   = {draw=ink, line width=0.55pt, fill=white, circle,
                  inner sep=0pt, minimum size=4.6mm, font=\scriptsize},
  asum/.style  = {draw=accent, line width=0.55pt, fill=white, circle,
                  inner sep=0pt, minimum size=4.6mm, font=\scriptsize},
  % 분기점
  dot/.style   = {ink, fill=ink, circle, inner sep=0pt, minimum size=1.5mm},
  % 신호 이름 — 선 위에 작게
  sname/.style = {font=\small, inner sep=1.5pt},
  % 그림 안의 짧은 설명. 그림 안의 글은 최소로 둔다
  note/.style  = {font=\scriptsize, soft, inner sep=1.5pt},
  anote/.style = {font=\scriptsize, accent, inner sep=1.5pt},
  % 축
  axis/.style  = {ink, line width=0.5pt},
  tick/.style  = {font=\scriptsize, soft},
}

% 캡션 한 줄 — 그림 아래 가는 선과 함께
\newcommand{\gnccaption}[2]{%
  \node[anchor=north west, text width=#1, align=left, font=\scriptsize, ink]
        at (current bounding box.south west) {\vspace{2pt}#2};}


\definecolor{bodyc}{HTML}{1D4ED8}
\definecolor{plantc}{HTML}{16A34A}

\begin{document}
\begin{tikzpicture}[x=1mm, y=1mm]

\pgfmathsetmacro{\psiD}{30}
\pgfmathsetmacro{\rot}{90-\psiD}
\pgfmathsetmacro{\SC}{12}                 % 1 m = 12 mm
\pgfmathsetmacro{\pbx}{2}                 % p^b = (2, 1) — 선체 축에서 읽은 값
\pgfmathsetmacro{\pby}{1}
\pgfmathsetmacro{\pnx}{round((\pbx*cos(\psiD) - \pby*sin(\psiD))*1000)/1000}
\pgfmathsetmacro{\pny}{round((\pbx*sin(\psiD) + \pby*cos(\psiD))*1000)/1000}
\pgfmathsetmacro{\SX}{\pny*\SC}           % 화면 x = 동 성분
\pgfmathsetmacro{\SY}{\pnx*\SC}           % 화면 y = 북 성분

% ===== 공통 — NED 축 =======================================================
\newcommand{\nedaxes}{%
  \draw[soft, line width=0.35pt, -{Stealth[length=3pt,width=2pt]}] (-3,0) -- (36,0);
  \draw[soft, line width=0.35pt, -{Stealth[length=3pt,width=2pt]}] (0,-3) -- (0,36);
  \node[font=\scriptsize, soft, anchor=north west] at (33,-0.6) {$y_n$ east};
  \node[font=\scriptsize, soft, anchor=south west] at (0.8,34)  {$x_n$ north};
}

% ===== 공통 — 선체와 선체 축 ===============================================
%   #1 선수축을 그릴 길이
\newcommand{\hull}[1]{%
  \begin{scope}[rotate=\rot]
    \draw[bodyc, line width=0.55pt, -{Stealth[length=3.4pt,width=2.4pt]}] (0,0) -- (#1,0);
    \node[font=\scriptsize, bodyc, anchor=west] at ({#1+1.2},0) {$x_b$ bow};
    \draw[bodyc, line width=0.55pt, -{Stealth[length=3.4pt,width=2.4pt]}] (0,0) -- (0,-15);
    \node[font=\scriptsize, bodyc, anchor=west] at (1.0,-15) {$y_b$ starboard};
    \draw[ink, line width=0.5pt, fill=white]
         (-4.2,-2.7) -- (-4.2,2.7) -- (1.6,2.7) -- (5.2,0) -- (1.6,-2.7) -- cycle;
  \end{scope}}

% ===== 공통 — 선수각 호 =====================================================
\newcommand{\psiarc}[1]{%
  \draw[errc, line width=0.42pt, -{Stealth[length=3pt,width=2pt]}]
       (0,#1) arc (90:\rot:#1mm);
  \pgfmathsetmacro{\amid}{(90+\rot)/2}
  \pgfmathsetmacro{\arad}{#1+3.8}
  \node[font=\scriptsize, errc] at ({\arad*cos(\amid)},{\arad*sin(\amid)}) {$\psi$};}

% ===========================================================================
% (1) 선체 축은 NED 축을 psi 만큼 돌린 것이다
% ===========================================================================
\begin{scope}[shift={(0,0)}]
  \node[font=\small\bfseries, anchor=north west, text width=64mm] at (-6,64)
       {(1) Two frames: one on the ground, one on the hull};
  \node[note, anchor=north west, text width=62mm] at (-6,55)
       {$x_n$ (north) and $y_n$ (east) are \textbf{fixed to the ground} and never move.
        $x_b$ (bow) and $y_b$ (starboard) are \textbf{fixed to the hull} and turn with it.
        The angle between them is the heading $\psi$};
  \nedaxes
  \hull{30}
  \psiarc{15}
  \node[anote, anchor=north west, text width=62mm] at (-6,-20)
       {The hull frame is the NED frame \textbf{turned} by $\psi$. That is the whole
        of what a rotation matrix does};
\end{scope}

% ===========================================================================
% (2) 벡터 하나를 두 축에서 읽는다
% ===========================================================================
\begin{scope}[shift={(76,0)}]
  \node[font=\small\bfseries, anchor=north west, text width=64mm] at (-6,64)
       {(2) One vector, two pairs of numbers};
  \node[note, anchor=north west, text width=62mm] at (-6,55)
       {The purple arrow $\mathbf{p}$ is a \textbf{single} vector: neither its length nor
        its direction changes. What changes is \textbf{which ruler reads it}};
  \nedaxes
  \hull{30}

  % 벡터 p — 주인공
  \draw[accent, line width=1.8pt, -{Stealth[length=5pt,width=4pt]}] (0,0) -- (\SX,\SY);
  \node[font=\small, accent, anchor=south west] at ({\SX+0.6},{\SY+0.6}) {$\mathbf{p}$};

  % NED 성분 — 가로 · 세로 점선
  \draw[accent, line width=0.35pt, densely dashed] (\SX,\SY) -- (\SX,0);
  \draw[accent, line width=0.35pt, densely dashed] (\SX,\SY) -- (0,\SY);
  \draw[accent, line width=1.0pt, -{Stealth[length=3.6pt,width=2.6pt]}] (0,0) -- (\SX,0);
  \draw[accent, line width=1.0pt, -{Stealth[length=3.6pt,width=2.6pt]}] (0,0) -- (0,\SY);
  %  동 성분 이름표는 축 **위**에 둔다. 축 아래에는 선체 성분의 점선이 지나가서
  %  글자와 겹친다 / the east-component label goes ABOVE the axis: below it the
  %  hull-component dashed line crosses the glyphs
  \node[font=\scriptsize, accent, anchor=north] at ({\SX*0.80},-1.2) {$y_n$};
  \node[font=\scriptsize, accent, anchor=east]  at (-1.2,{\SY*0.52}) {$x_n$};

  % 선체 성분 — 돌아간 축을 따라
  \begin{scope}[rotate=\rot]
    \draw[errc, line width=0.35pt, densely dashed]
         ({\pbx*\SC},{-\pby*\SC}) -- ({\pbx*\SC},0);
    \draw[errc, line width=0.35pt, densely dashed]
         ({\pbx*\SC},{-\pby*\SC}) -- (0,{-\pby*\SC});
    \draw[errc, line width=1.0pt, -{Stealth[length=3.6pt,width=2.6pt]}]
         (0,0) -- ({\pbx*\SC},0);
    \draw[errc, line width=1.0pt, -{Stealth[length=3.6pt,width=2.6pt]}]
         (0,0) -- (0,{-\pby*\SC});
  \end{scope}
  %  두 성분의 이름표는 **돌리지 않은** 화면 좌표에 둔다. 돌린 좌표에 두면 원점의
  %  선체 그림 위로 올라앉는다 / the two component labels are placed in unrotated
  %  screen coordinates: inside the rotated scope they land on top of the hull
  \node[font=\scriptsize, errc, anchor=east] at (6.2,16.0) {$x_b$};
  \node[font=\scriptsize, errc, anchor=north] at (5.2,-7.6) {$y_b$};

  \node[font=\scriptsize, errc, anchor=north west] at (-6,-20)
       {read with the hull ruler\; $\mathbf{p}^{b} = (2.000,\; 1.000)$};
  \node[font=\scriptsize, accent, anchor=north west] at (-6,-26)
       {read with the ground ruler\;\; $\mathbf{p}^{n} = (\pgfmathprintnumber[fixed,precision=3]{\pnx},\;
        \pgfmathprintnumber[fixed,precision=3]{\pny})$};
\end{scope}

% ===========================================================================
% (3) R 의 열은 선체 축을 NED 에서 읽은 것이다
% ===========================================================================
\begin{scope}[shift={(152,0)}]
  \node[font=\small\bfseries, anchor=north west, text width=64mm] at (-6,64)
       {(3) The columns of $R$ are nothing to memorise};
  \node[note, anchor=north west, text width=62mm] at (-6,55)
       {A column of $R$ is one \textbf{unit vector} of the hull frame, read with the
        ground ruler. The first column is $\mathbf{e}_{x_b}$, the second is
        $\mathbf{e}_{y_b}$};
  \nedaxes
  \hull{30}

  % 선체 x 축 단위벡터  (NED 성분 = (cos psi, sin psi))
  \pgfmathsetmacro{\ux}{round(sin(\psiD)*1000)/1000}     % 동 성분
  \pgfmathsetmacro{\uy}{round(cos(\psiD)*1000)/1000}     % 북 성분
  \draw[alng, line width=1.6pt, -{Stealth[length=4.6pt,width=3.6pt]}]
       (0,0) -- ({\ux*22},{\uy*22});
  \draw[alng, line width=0.35pt, densely dashed] ({\ux*22},{\uy*22}) -- ({\ux*22},0);
  \draw[alng, line width=0.35pt, densely dashed] ({\ux*22},{\uy*22}) -- (0,{\uy*22});
  \node[font=\scriptsize, alng, anchor=south east] at ({\ux*22-0.8},{\uy*22+1.6}) {$\mathbf{e}_{x_b}$};

  % 선체 y 축 단위벡터  (NED 성분 = (-sin psi, cos psi))
  \pgfmathsetmacro{\vx}{round(cos(\psiD)*1000)/1000}     % 동 성분
  \pgfmathsetmacro{\vy}{round(-sin(\psiD)*1000)/1000}    % 북 성분
  \draw[plantc, line width=1.6pt, -{Stealth[length=4.6pt,width=3.6pt]}]
       (0,0) -- ({\vx*22},{\vy*22});
  \node[font=\scriptsize, plantc, anchor=north] at ({\vx*22},{\vy*22-1.2}) {$\mathbf{e}_{y_b}$};

  \node[font=\scriptsize, alng, anchor=north west] at (-6,-20)
       {$\mathbf{e}_{x_b} = (\cos\psi,\;\sin\psi) = (0.866,\; 0.500)$};
  \node[font=\scriptsize, plantc, anchor=north west] at (-6,-26)
       {$\mathbf{e}_{y_b} = (-\sin\psi,\;\cos\psi) = (-0.500,\; 0.866)$};
\end{scope}

% ===========================================================================
% 식 — 아래 띠
% ===========================================================================
\draw[wire, soft] (-6,-36) -- (216,-36);

\node[anchor=north west, font=\small] at (-6,-41)
     {$\underbrace{%
        \begin{bmatrix} x_n \\ y_n \end{bmatrix}}_{\text{in NED}}
       =
       \underbrace{%
        \begin{bmatrix} \cos\psi & -\sin\psi \\ \sin\psi & \phantom{-}\cos\psi \end{bmatrix}}_{R(\psi)}
       \underbrace{%
        \begin{bmatrix} x_b \\ y_b \end{bmatrix}}_{\text{in the hull}}$};

\node[anchor=north west, font=\small] at (66,-41)
     {$\begin{bmatrix} x_b \\ y_b \end{bmatrix}
       = R(\psi)^{\mathsf T}\begin{bmatrix} x_n \\ y_n \end{bmatrix}
       = \begin{bmatrix} \phantom{-}\cos\psi & \sin\psi \\ -\sin\psi & \cos\psi \end{bmatrix}
         \begin{bmatrix} x_n \\ y_n \end{bmatrix}$};

\node[anchor=north west, font=\scriptsize, warn, text width=66mm] at (152,-40)
     {\textbf{Reading the other way is a transpose.} $R^{-1} = R^{\mathsf T}$,\;
      $R R^{\mathsf T} = I$,\; $\det R = 1$.
      No inverse is ever computed: the rows and the columns are swapped};

\node[anchor=north west, font=\scriptsize, soft, text width=215mm] at (-6,-62)
     {Numerical check at $\psi = 30^{\circ}$:\;
      $R = \begin{bmatrix} 0.866 & -0.500 \\ 0.500 & \phantom{-}0.866 \end{bmatrix}$,\quad
      each column has length 1,\quad the two columns have inner product 0,\quad $\det R = 1$.\quad
      In MATLAB this is one line, \texttt{R = [cos(psi) -sin(psi); sin(psi) cos(psi)]}};

\end{tikzpicture}
\end{document}
